Probatutako konfigurazioen artean, P kontrolagailuak $K_p = 1{,}5$ balioarekin eta $0{,}2\text{ m/s}$-ko abiadura konstantearekin eman du emaitzarik orekatuena. Robotak $3{,}32$ buelta egin ditu talkarik gabe, zuzenetako $\theta_{\text{RMS}}$ $1{,}7^\circ$ izan da eta paretarekiko batez besteko distantzia $0{,}307\text{ m}$ izan da. $K_p = 2{,}0$ balioarekin ere $3{,}32$ buelta egin ditu talkarik gabe, eta zuzenetako $\theta_{\text{RMS}}$ berdina izan da. Hala ere, $K_p = 1{,}5$ aukeratu dugu, paretarekiko distantzia egonkorragoa izan delako zuzenetan, desbideratze estandarra $0{,}013\text{ m}$ izan baita, beste konfigurazioko $0{,}028\text{ m}$-ren aldean.

Probek erakutsi dute irabazi txikiekin robotak berandu biratzen duela, eta irabazi handiekin, berriz, gehiegi zuzentzen duela norabidea eta pareta galtzen duela. Kurbetan moteltzea lagungarria izan da $K_p = 1{,}0$ balioarekin, baina $K_p = 1{,}5$ eta $2{,}0$ balioekin biraketa gehiegi itxi du. Horregatik, aukeratutako $K_p = 1{,}5$ balioarekin abiadura konstantea mantendu dugu.

I osagaiak kurbetako errorea noranzko berean pilatu du. Horren ondorioz, kurba amaitu ondoren ere robotak biratzen jarraitu du eta zuzenetan paretatik gehiago urrundu da. D osagai txikiek ez dute P kontrolaren emaitza hobetu, eta $K_d = 1{,}0$ balioak biraketa bortitzak sortu ditu. $K_p = 2{,}0$ finkatuta, P kontrolak eta probatutako PID konbinazioak ibilbidea talkarik gabe osatu dute, eta zuzenetako angelu-errorea antzekoa izan da. Proba horretan I eta D gehitzeak hobekuntza argirik ekarri ez duenez, P kontrolarekin geratu gara.
